\section{When Presence Is Not Felt}

The most ordinary obstacle to prayer is not an argument but an absence: nothing seems to answer. God feels distant; prayer feels inert; suffering looks like neglect or punishment. The mind moves quickly from an experience to a verdict: God is not there; or he is there but does not care; or he is making me suffer because I have failed.

None of those conclusions follows from the feeling.

The first part of this series argued that the creature exists only as continuously received. If God ceased to cause a creature's act of being, there would be no abandoned thing left over to register the abandonment. At the instant in which prayer seems to strike a ceiling, the person praying is being held in existence. This ontological presence is not the same as sanctifying indwelling, sacramental communion, or sensible consolation; those must not be collapsed. But sensible dryness by itself proves neither divine absence nor loss of grace.

Sin can wound friendship with God, and known grave sin calls for repentance and sacramental Reconciliation. Yet a bare feeling of distance is not a moral diagnosis. It may accompany spiritual purification; it may also accompany grief, exhaustion, illness, distraction, depression, ordinary change in affect, or no cause that can be confidently named. To label every dryness a ``dark night'' can flatter the sufferer and presume to read God's particular action. To label it punishment is no safer.

St.~Teresa of Jesus calls contemplative prayer a ``close sharing between friends'' and the taking of time to be alone with the One who loves us. The definition is spare enough to survive an hour in which the sharing is not felt. The \emph{Catechism} immediately adds that one makes time with a determined will and does not give up because of trials or dryness.\endnote{CCC 2709--2719, especially 2709--2712; Teresa of Jesus, \emph{The Book of Her Life} 8.5.}

John of the Cross can describe passive purification with terrifying accuracy, but his doctrine is a map for discernment, not a label for every barren Tuesday. Its controlling end is union of love; its agent is God; its ordinary measure is the growth of faith, hope, charity, detachment from disorder, and concrete fidelity. The tradition does not authorize a person to neglect natural remedies or pastoral counsel because a spiritual explanation sounds more elevated.\endnote{John of the Cross, \emph{Dark Night} I.8--10 and II.5--8; Benedict XVI, general audience, 16 February 2011.}

Thérèse of Lisieux gives a different witness. In her final trial against faith, darkness did not become evidence that prayer was false. She renewed acts of faith and interceded for unbelievers from within the darkness. Her witness does not romanticize desolation; it shows love refusing to make sensible assurance the price of fidelity.\endnote{Francis, \emph{C'est la confiance} (15 October 2023), nos.~25--27, receiving Thérèse's public writings.}

Prayer is therefore not a technique for generating the experience of presence. It is an enacted truth: the creature turns toward the One in whom it already exists, whom faith knows in Christ, and whom charity desires even when the emotions cannot follow. Dryness can make this act poor. It cannot make it unreal.
